\subsection{章末確認問題}
\begin{enumerate}[leftmargin=2em]
  \item ソフトウェアエンジニアリング管理は、単なる進捗管理ではない。理由を説明せよ。
  \item ソフトウェアが物理製品と異なるために、管理が難しくなる点を三つ挙げよ。
  \item 管理と測定を切り離すと、どのような問題が起きるかを説明せよ。
  \item 開始と範囲定義で、要求の決定と交渉が重要な理由を説明せよ。
  \item 実現可能性分析で確認すべき観点を五つ挙げよ。
  \item 予測型と適応型のライフサイクルの違いを、要求、計画、利害関係者関与の観点で説明せよ。
  \item 工数・日程・費用見積りがソフトウェアで難しい理由を説明せよ。
  \item リスク管理で行う識別、分析、対応、監視の違いを説明せよ。
  \item 品質管理をテストだけに任せてはいけない理由を説明せよ。
  \item ソフトウェア取得で、契約に明記すべき事項を挙げよ。
  \item 監視プロセスと制御プロセスの違いを説明せよ。
  \item 終結活動で、教訓を組織の学習へつなげることが重要な理由を説明せよ。
  \item 情報ニーズ、測定量、情報製品の関係を説明せよ。
\end{enumerate}
\begin{center}
\fbox{\begin{minipage}{0.94\linewidth}
解答の方向性\par\smallskip
1 は、計画、見積り、測定、制御、調整、リスク管理、リーダーシップを含む点に注目する。2 は、要求変更、無形性、可鍛性、技術変化、人間要因、継続的改善を挙げる。3 は、根拠のない判断と目的のない数値集めを対比する。4 以降は、本文の「合意」「実現可能性」「範囲」「測定」「監視と制御」「学習と改善」という語を使って説明できる。
\end{minipage}}
\end{center}
